% Problem record op_f271040c57d9a513. This TeX file is derived from
% database/problems_json/op_f271040c57d9a513.json by scripts/sync-tex.mjs; edit the JSON record
% and rerun the script rather than editing this file.
\section{Scalable pseudorandom unitaries with an independent security parameter}
\label{sec:op_f271040c57d9a513}

\paragraph{Problem.}
Can pseudorandom-unitary security scale independently of Hilbert-space dimension while construction uses only polynomially many oracle queries?

Question. Can pseudorandom-unitary security be increased independently of Hilbert-space dimension while retaining polynomial construction-query cost?

Use the precise idealized model of \sourcecite{ref:f271-brakerski26}{Brakerski26}, Definitions 4.1 and 4.4. For every dimension $d=2^n$ and security parameter $\kappa$, seek a uniformly generated family $\{U_{n,\kappa,f}\}_f$, indexed by Boolean functions $f:\{0,1\}^{r}\to\{0,1\}$, with

\begin{equation}
r,\ q\leq\operatorname{poly}(n,\kappa).
\label{eq:f271-1}
\end{equation}

A construction algorithm $\mathcal A^f$ uses at most $q$ oracle queries and satisfies, for every $f$,

\begin{equation}
\left\|\mathcal A^f-
U_{n,\kappa,f}(\,\cdot\,)U_{n,\kappa,f}^{\dagger}
\right\|_{\diamond}\leq2^{-\kappa}.
\label{eq:f271-2}
\end{equation}

For uniformly random $f$, require

\begin{equation}
\sup_{\mathcal D:\,\#\mathrm{queries}\leq2^{\kappa}}
\left|\Pr_f[\mathcal D^{U_{n,\kappa,f}}=1]
-\Pr_{V\sim\operatorname{Haar}(d)}[\mathcal D^V=1]\right|
\leq2^{-\kappa}.
\label{eq:f271-3}
\end{equation}

The distinguisher accesses the ideal unitary, not the underlying function $f$. Both algorithms are query-bounded, not necessarily computationally bounded. There must be no imposed scaling relation between $n$ and $\kappa$. \sourcecite{ref:f271-brakerski26}{Brakerski26}

The displayed definitions, constraints, and target bounds are recorded in Eqs.~\eqref{eq:f271-1}, \eqref{eq:f271-2}, \eqref{eq:f271-3}.

\subsection*{Status}

Unsolved

\subsection*{Source}

The question is explicitly posed or retained as open in the cited primary literature \sourcecite{ref:f271-brakerski26}{Brakerski26}.  The statement is rewritten here to make its hypotheses and success criterion self-contained.

\subsection*{Progress}

\begin{itemize}
  \item Status: explicitly unresolved in the 2026 source. Brakerski and Yuen analyze this scalable formulation, establish barriers for existing candidates, and connect a positive answer to unitary synthesis. Their baseline construction has polynomial dependence on $d$, rather than the required $\log d$. Thus it does not resolve the displayed problem. \sourcecite{ref:f271-brakerski26}{Brakerski26}

  \item Their distinguishing attack on the PFC construction is not a universal impossibility theorem for scalable PRUs, and the paper leaves the general construction problem unresolved. \sourcecite{ref:f271-brakerski26}{Brakerski26}
\end{itemize}

\subsection*{References}

\begin{enumerate}[label={},leftmargin=4.8em,labelsep=0.5em]
  \footnotesize
  \item[\textup{[Brakerski26]}]\label{ref:f271-brakerski26}
  Zvika Brakerski and Henry Yuen, \emph{On Scalable Pseudorandom Unitaries and the Unitary Synthesis Problem}. \href{https://arxiv.org/html/2605.09957v1}{arXiv:2605.09957v1}, submitted 11 May 2026; CRYPTO 2026. See Definitions 4.1 and 4.4, §4.1, Lemma 4.7, and §6.
\end{enumerate}

\subsection*{Comment}

This is deliberately an oracle-query problem. A query-efficient construction is not automatically a gate-efficient, standard-model cryptosystem. That distinction should remain explicit in any proposed solution.

The independence of parameters is substantive. A distinguishing bound involving $t^2/2^n$ may become small by increasing $n$; it does not offer an independent security knob at fixed dimension. The requested construction must work when the allowed query count is large relative to dimension as well.

The object is a quantum operation, and the security criterion is operational distinguishability from Haar rather than a relationship between classical complexity classes.

\subsection*{Field}

Quantum Cryptography

\subsection*{Topic}

Computational complexity and computability; Quantum circuit complexity

\subsection*{ID}

\texttt{op\_f271040c57d9a513}
